\documentclass[12pt]{article}

\textwidth 16cm \textheight 23cm \evensidemargin 0cm
\oddsidemargin 0cm \topmargin -2cm
\parindent 0pt
\parskip \medskipamount


\usepackage[dutch]{babel}
\usepackage{amssymb}
\usepackage{amsmath}
\usepackage[utf8,latin1]{inputenc}
\usepackage[normalem]{ulem} % strikethrough normal text with \sout{text}
\usepackage{cancel} % strikethrough in math mode with \cancel{text}
%\usepackage{gensymb}

\usepackage{subfig}
\usepackage{wrapfig}
\usepackage{graphicx}
\usepackage{graphics}
\usepackage{latexsym}
%\usepackage[latin1]{inputenc} %letters met accenten mogen in broncode gebruikt worden
\usepackage{amssymb,amsmath,amsthm} %allerlei wiskundige symbolen => http://www.ams.org/tex/amslatex.html voor read-me's
%\usepackage[small,bf,hang]{caption}
\usepackage[table]{xcolor}
\usepackage{pgf,tikz}
\usetikzlibrary{arrows}

\usepackage[pdftex]{hyperref} 
\hypersetup{setpagesize=false}
%\hypersetup{bookmarksopen}
\hypersetup{bookmarksnumbered} %voor de bookmarks in pdf de nrs zetten ook
\hypersetup{pdfborder=0.1pt} %de rand rond de links smaller maken
\hypersetup{linktocpage=true} %titels of pgnrs omkaderen in toc
\hypersetup{colorlinks=true} %kaders rond links of tekst in kleur, kaders gaan weg bij afprinten, kleur blijft


\numberwithin{equation}{section} %om de equations labels te geven die met section meegaan, dus niet gwn (1) (2) ... maar (1.1) (1.2) ...

%iets voor modulonotatie
\makeatletter
\def\imod#1{\allowbreak\mkern10mu({\operator@font mod}\,\,#1)}
\makeatother

%Getallenverzamelingen
\newcommand{\N}{\mathbb{N}}
\newcommand{\Z}{\mathbb{Z}}
\newcommand{\Q}{\mathbb{Q}}
\newcommand{\R}{\mathbb{R}}
\newcommand{\C}{\mathbb{C}}
\newcommand{\mO}{\mathcal{O}}
%<= en >= mooier maken
\renewcommand{\leq}{\leqslant}
\renewcommand{\geq}{\geqslant}

%speciale wiskunde letters: \mathcal{} of \mathfrak{}

%Nieuwe omgevingen voor definities en stellingen
\theoremstyle{plain}  \newtheorem{stel}{Stelling}[section]
\theoremstyle{plain}  \newtheorem{lemma}[stel]{Lemma}
\theoremstyle{plain}  \newtheorem{gevolg}[stel]{Gevolg}
\theoremstyle{definition}  \newtheorem{defi}[stel]{Definitie}
\theoremstyle{definition} \newtheorem{vb}[stel]{Voorbeeld}
\theoremstyle{definition} \newtheorem{vbn}[stel]{Voorbeelden}
\theoremstyle{definition}  \newtheorem{opm}[stel]{Opmerking}
\theoremstyle{plain}  \newtheorem{opdracht}{Opdracht}

\hyphenation{e-ner-zijds}
\hyphenation{ka-nons-ko-gel}
\hyphenation{schrij-ven}
\hyphenation{ka-nons-ko-gels}
\hyphenation{ont-stond}
\hyphenation{kwa-dra-tisch}

\usepackage{color}
\newcommand{\todo}[1]{\textcolor{red}{\##1\#}}
\newcommand{\question}[1]{\textcolor{blue}{\##1\#}}

% Wissel tussen versie van leerkracht en versie voor leerlingen
% gebruik: \teacher{Dit is een notitie voor de leerkracht}
% eerste parameter is de notitie
\newcommand{\teacher}[1]{\textcolor{gray}{\emph{#1}}}
%\newcommand{\teacher}[1]{}

% Wissel tussen versie van leerkracht en versie voor leerling door antwoorden te verbergen
% gebruik: \answer[2cm]{Dit is een antwoord}
% tussen vierkante haakjes is optioneel de lengte die voorzien moet worden (default = 0cm),
% de parameter is het antwoord
\newcommand{\answer}[2][0cm]{{\textcolor{gray}{\emph{#2}}}}
%\newcommand{\answer}[2][0cm]{\vspace*{#1}}

\graphicspath{{figuren/}}

%\newtheorem{definitie}{Definitie}

\newcommand{\degree}{\ensuremath{^\circ}}

\begin{document}

\thispagestyle{empty}

\begin{center}
 \begin{minipage}{7cm}
\begin{center}
%\aula 
%B


\includegraphics[width=5cm]{ruglogo.pdf}


 {\large
Universiteit Gent \\
Faculteit  Wetenschappen \\
Vakgroep  Wiskunde
} \end{center}\end{minipage}

\vspace{3cm}

 {\Huge\bf 
  Vullen van Vlakken en Volumes:\\ Administratie}
  
\baselineskip=12pt
\vspace{1.5cm}

{\large\bf  Lien Lambert }\\
\vspace{0.5cm}
{\large \bf Pieter Pareit }\\
\vspace{0.5cm}
{\large \bf Jordy Vanpoucke}\\

\baselineskip=12pt

\vspace{2cm}

 
\bigskip

{\large Academiejaar 2011-2012}

%\vspace{5cm}
%\baselineskip=1.2cm



\end{center}

\vfill

\hspace*{\fill}
\begin{minipage}{8cm}
\noindent   Project Vakdidactiek Wiskunde II 
 \end{minipage}


\bigskip

\hspace*{\fill}
\begin{minipage}{8cm}
\noindent  Prof. Dr.  F. De Clerck\\
Praktijkassistent G. Maes
 \end{minipage}

\newpage
\tableofcontents 
\newpage

\section{Doelgroep}
Leerlingen uit een ASO-richting met de pool wiskunde.

\section{Voorkennis}
In dit lessenpakket over het vullen van vlakken en volumes steunen we expliciet op de geziene leerstof uit
\begin{itemize}
\item de getallenleer en algebra: De leerlingen
\begin{itemize}
\item (19) kunnen vergelijkingen van de eerste en de tweede graad in \'{e}\'{e}n onbekende oplossen,
\item (20) kunnen ongelijkheden van de eerste en de tweede graad in \'{e}\'{e}n onbekende oplossen,
\item (21) lossen problemen op die kunnen vertaald worden naar:
\begin{itemize}
\item een vergelijking van de eerste en de tweede graad in \'{e}\'{e}n onbekende.
\item een ongelijkheid van de eerste en de tweede graad in \'{e}\'{e}n onbekende.
\end{itemize}
\end{itemize}
\item de re\"{e}le functies: De leerlingen
\begin{itemize}
\item (28) kunnen stelsels van twee vergelijkingen van de eerste graad met twee onbekenden algebra�sch oplossen en de oplossing grafisch interpreteren,
\item (29) kunnen problemen oplossen die te vertalen zijn naar stelsels van twee vergelijkingen van de eerste graad met twee onbekenden,
\item (31) lossen problemen op die kunnen beschreven worden met eerste- en tweedegraadsfuncties.
\end{itemize}
\item de meetkunde: De leerlingen
\begin{itemize}
\item (34) verklaren gelijkvormigheid van figuren met behulp van schaal en congruentie,
\item (35) gebruiken de gelijkvormigheid van driehoeken en de stelling van Thales om de lengte van lijnstukken te berekenen,
\item (39) kunnen problemen met zijden en hoeken van driehoeken uit de technische wereld oplossen door een effici�nte keuze te maken uit:
\begin{itemize}
\item de stelling van Thales
\item de stelling van Pythagoras
\item goniometrische getallen
\end{itemize}
\item (41) lossen eenvoudige problemen i.v.m. ruimtelijke situaties op door gebruik te maken van eigenschappen van vlakke figuren,
\item (43) kunnen de inhoud van sommige ruimtelijke objecten benaderend berekenen door ze op te splitsen in of aan te vullen tot gekende lichamen,
\item (44) kunnen effecten van schaalverandering op inhoud en oppervlakte berekenen.
\end{itemize}
\end{itemize}


\section{Globale motivatie}
Voor deze lessenreeks zijn we vertrokken vanuit het bolstapelprobleem van Kepler. Het viel ons op hoe men in het dagelijkse leven bezig is met het zo optimaal mogelijk stapel van bollen, maar ook van kubussen en dergelijke. In deze lessenreeks wordt dit opgebouwd door te starten met het tweedimensionale probleem, namelijk de vlakvulling. In deze lessenreeks is het de bedoeling dat de leerlingen zelf gaan experimenteren, en hun bevindingen ook hard maken via het bewijzen of uitrekenen ervan.

\section{Eindtermen}

\subsection{Vakoverschrijdende eindtermen}
In het lessenpakket komen de volgende vakoverschrijdende eindtermen aan bod.
\subsubsection{Gemeenschappelijke stam}
De leerlingen	 
\begin{itemize}
\item	(2) kunnen originele idee\"{e}n en oplossingen ontwikkelen en uitvoeren,
\item (4)	blijven, ondanks moeilijkheden, een doel nastreven,
\item (6)	kunnen schoonheid ervaren, 
\item (7)	kunnen schoonheid cre\"{e}ren,
\item (8)	benutten leerkansen in diverse situaties,
\item (10) engageren zich spontaan,
\item (11) kunnen gegevens, handelwijzen en redeneringen ter discussie stellen a.d.h. van relevante criteria,
\item (12) zijn bekwaam om alternatieven af te wegen en een bewuste keuze te maken,
\item (13) kunnen onderwerpen benaderen vanuit verschillende invalshoeken,
\item (19) dragen actief bij tot het realiseren van gemeenschappelijke doelen.
\end{itemize}
\subsubsection{Leren leren}
De leerlingen
\begin{itemize}
\item (4)	kunnen verwerkte informatie vakoverstijgend en in verschillende situaties functioneel toepassen,
\item (5)	kunnen informatie samenvatten,
\item (6)	kunnen op basis van hypothesen en verwachtingen mogelijke oplossingswijzen realistisch inschatten en uitvoeren,
\item (7)	evalueren de gekozen oplossingswijze en de oplossing en gaan eventueel op zoek naar een alternatief.
\end{itemize}
\subsubsection{Contexten}
We maken enkel gebruik van eindtermen uit context 7, `socioculturele samenleving'.\\
De leerlingen
\begin{itemize}
\item (6) gaan actief om met de cultuur en kunst die hen omringen,
\item (7)	illustreren de wederzijdse be�nvloeding van kunst, cultuur en techniek, politiek, economie, wetenschappen en levensbeschouwing.
\end{itemize}
\subsection{Vakgebonden eindtermen}
We maken gebruik van de volgende vakgebonden eindtermen.\\
De leerlingen kunnen
\begin{itemize}
\item (1)	wiskundetaal begrijpen en gebruiken,
\item (2)	wiskundige informatie analyseren, schematiseren en structureren,
\item (3)	eenvoudig mathematiseerbare problemen ontleden (onderscheid maken tussen gegevens en gevraagde, de relevantie van de gegevens nagaan en verbanden leggen ertussen) en vertalen naar een passende wiskundige context,
\item (4)	wiskundige problemen planmatig aanpakken (door eventueel hi\"{e}rarchisch op te splitsen in deelproblemen),
\item (5)	bij het oplossen van wiskundige problemen kritisch reflecteren over het oplossingsproces en het eindresultaat,
\item (8)	voorbeelden geven van de rol van de wiskunde in de kunst,
\item (9)	kennis, inzicht en vaardigheden die ze verwerven in wiskunde bij het verkennen, vertolken en verklaren van problemen uit de realiteit gebruiken.
\end{itemize}
De leerlingen
\begin{itemize}
\item (11*)	leggen een zin voor nauwkeurigheid aan de dag bij het hanteren en het toepassen van wiskunde,
\item (13*)	zijn gericht op samenwerking om de eigen mogelijkheden te vergroten.
\end{itemize}
\subsection{Decretale specifieke eindtermen wiskunde}
Bij dit lessenpakket horen de volgende decretale specifieke eindtermen.\\
De decretale specifieke eindtermen wiskunde hebben betrekking op kennis, inzichten, vaardigheden en attitudes waarmee leerlingen:
\begin{itemize}
\item verbanden leggen tussen wiskunde en praktische toepassingen uit het dagelijkse leven en zo relaties leggen met problemen uit maatschappij, wetenschap en techniek,
\item een wiskundig denken en redeneren ontwikkelen, d.w.z. een wiskundig eigen wijze van: 
	\begin{itemize}
	\item bevragen, onderzoeken en formuleren van vermoedens,
	\item modelleren en structureren,
	\item argumenteren en bewijzen,
	\end{itemize}
\item gesloten en open problemen wiskundig kunnen stellen en analyseren, en oplossingen argumenteren en bespreken,
\item communiceren over wiskundig beschreven situaties, met inbegrip van het vlotte gebruik van meer specifieke wiskundetaal,
\item kritisch reflecteren op denken en handelen.
\end{itemize}

\section{Planning}
\subsection{Les 1: Inleiding}


\subsection{Les + : Tweedimensionale problemen}

\subsubsection{Algemeen en zo}
De leerlingen gaan in groepjes zitten aan een bank en overlopen samen met de leerkracht het bundeltje. Het is vooral belangrijk om stil te staan bij de 'vragen' en 'opdrachten', die ze dan per groepje mogen proberen oplossen.

\subsubsection{Experimenteren}



\section{Referenties}
\begin{thebibliography}{1}
\bibitem{1}\label{1} Hans Meissen en Rob van Oord, \textit{Schuiven met auto's, munten en bollen}, Epsilon Uitgaven, Utrecht, 2001-2002\\
\end{thebibliography}
\end{document}